\let\negmedspace\undefined
\let\negthickspace\undefined
\documentclass[journal]{IEEEtran}
\usepackage[a5paper, margin=10mm, onecolumn]{geometry}
%\usepackage{lmodern} % Ensure lmodern is loaded for pdflatex
\usepackage{tfrupee} % Include tfrupee package

\setlength{\headheight}{1cm} % Set the height of the header box
\setlength{\headsep}{0mm}     % Set the distance between the header box and the top of the text

\usepackage{gvv-book}
\usepackage{gvv}
\usepackage{cite}
\usepackage{amsmath,amssymb,amsfonts,amsthm}
\usepackage{algorithmic}
\usepackage{graphicx}
\usepackage{textcomp}
\usepackage{xcolor}
\usepackage{txfonts}
\usepackage{listings}
\usepackage{enumitem}
\usepackage{mathtools}
\usepackage{gensymb}
\usepackage{comment}
\usepackage[breaklinks=true]{hyperref}
\usepackage{tkz-euclide} 
\usepackage{listings}
% \usepackage{gvv}                                        
\def\inputGnumericTable{}                                 
\usepackage[latin1]{inputenc}                                
\usepackage{color}                                            
\usepackage{array}                                            
\usepackage{longtable}                                       
\usepackage{calc}                                             
\usepackage{multirow}                                         
\usepackage{hhline}                                           
\usepackage{ifthen}                                           
\usepackage{lscape}
\begin{document}

\bibliographystyle{IEEEtran}
\vspace{3cm}

\title{12.359}
\author{EE25BTECH11033 - Kavin}
% \maketitle
% \newpage
% \bigskip
{\let\newpage\relax\maketitle}

\renewcommand{\thefigure}{\theenumi}
\renewcommand{\thetable}{\theenumi}
\setlength{\intextsep}{10pt} % Space between text and floats
\textbf{Question}:\\
Which one of the following statements is NOT CORRECT?

\begin{enumerate}
\item A well-conditioned matrix has a condition number close to 1  
\item An ill-conditioned matrix has a large condition number  
\item The inverse of a well-conditioned matrix can be computed accurately  
\item A non-invertible matrix has a condition number close to 1  
\end{enumerate}\\    
\bigskip

\textbf{Solution}:\\
The \textbf{condition number} of an invertible matrix $\vec{A}$, denoted as $\kappa(\vec{A})$, is defined as:
\begin{align}
    \kappa(\vec{A}) = \norm{\vec{A}} \norm{\vec{A}^{-1}}
\end{align}
where,
\begin{align*}
    \norm{\vec{A}} = \text{norm of matrix $\vec{A}$}
\end{align*}
It measures how much the output of a linear system $\vec{A}\vec{x}=\vec{b}$ can change for a small change in the input data.\\
\\
The most theoretically important and widely used condition number is the one defined with respect to the \textbf{matrix 2-norm} (also known as the spectral norm).

The 2-norm of a matrix $\vec{A}$, $\norm{\vec{A}}_2$, is its largest singular value, $\sigma_{\text{max}}(\vec{A})$. The 2-norm of its inverse, $\norm{\vec{A^{-1}}}_2$, is the reciprocal of the smallest singular value of $\vec{A}$, $1/\sigma_{\text{min}}(\vec{A})$.

Therefore, the \textbf{spectral condition number}, denoted $\kappa_2(\vec{A})$, is defined as the ratio of the largest to the smallest singular value of the matrix:
\begin{align}
\kappa_2(\vec{A}) = \frac{\sigma_{\text{max}}(\vec{A})}{\sigma_{\text{min}}(\vec{A})}
\end{align}
where:
\begin{itemize}
    \item $\sigma_{\text{max}}(\vec{A})$ is the largest singular value of $\vec{A}$.
    \item $\sigma_{\text{min}}(\vec{A})$ is the smallest singular value of $\vec{A}$.
\end{itemize}

\textbf{For a symmetric positive definite matrix}, the singular values are the same as the eigenvalues. In this special case, the definition simplifies to the ratio of the largest to the smallest eigenvalue:
\begin{align}
\kappa_2(\vec{A}) = \frac{\lambda_{\text{max}}(\vec{A})}{\lambda_{\text{min}}(\vec{A})} \quad (\text{for symmetric positive definite } \vec{A})
\end{align}


\section*{What the Condition Number Measures}

The condition number provides a measure of the sensitivity of the solution of a linear system of equations $\vec{A}\vec{x} = \vec{b}$ to errors or perturbations in the input data ($\vec{A}$ or $\vec{b}$).

Specifically, it bounds the maximum possible relative error in the solution $x$ for a given relative error in the input vector $b$. The fundamental inequality is:
\begin{align}
\frac{\norm{\delta \vec{x}}}{\vec{x}} \le \kappa(\vec{A}) \frac{\norm{\vec{b}}}{\vec{b}}
\end{align}

\section*{Interpretation}
\begin{itemize}
    \item \textbf{If $\kappa(\vec{A})$ is small (close to 1)}, the matrix is \textbf{well-conditioned}. This means that small relative errors in $\vec{b}$ will only lead to small relative errors in the solution $\vec{x}$. The solution is stable and reliable.
    \item \textbf{If $\kappa(\vec{A})$ is large}, the matrix is \textbf{ill-conditioned} or \textbf{poorly-conditioned}. This means that a tiny relative error in $\vec{b}$ can be magnified into a very large relative error in the solution $\vec{x}$. The solution is unstable and may be unreliable from a numerical perspective.
\end{itemize}\\
\bigskip


\begin{itemize}
    \item[\textbf{1)}] \textbf{A well-conditioned matrix has a condition number close to 1:} This is \textbf{CORRECT}. The smallest possible condition number is 1 (for example, the identity matrix). A value close to 1 indicates that the matrix is numerically stable.\\

    \item[\textbf{2)}] \textbf{An ill-conditioned matrix has a large condition number:} This is \textbf{CORRECT}. This is the definition of an ill-conditioned matrix. A large condition number means that small errors in the input can lead to very large errors in the solution.\\

    \item[\textbf{3)}] \textbf{The inverse of a well-conditioned matrix can be computed accurately:} This is \textbf{CORRECT}. Because a well-conditioned matrix is not sensitive to small perturbations, numerical algorithms used to compute its inverse will produce results with high accuracy.\\

    \item[\textbf{4)}] \textbf{A non-invertible matrix has a condition number close to 1:} This is \textbf{NOT CORRECT}. A non-invertible (or singular) matrix does not have an inverse $\vec{A}^{-1}$. For the purposes of conditioning, its inverse is considered to have an infinite norm. Therefore, the condition number of a non-invertible matrix is defined to be \textbf{infinity ($\infty$)}. An infinite condition number is the largest possible value, representing the worst possible conditioning, not the best (which is close to 1).
\end{itemize}


\end{document}


